\documentclass[11pt]{article}

% --- Packages ---------------------------------------------------------------
\usepackage[a4paper,margin=1in]{geometry}
\usepackage{amsmath,amssymb}
\usepackage{booktabs}
\usepackage{array}
\usepackage{microtype}
\usepackage{titlesec}
\usepackage{xcolor}
\usepackage{fancyhdr}
\usepackage[hidelinks]{hyperref}

% --- Bordeaux color (sampled from the Codex Vini house style) ---------------
\definecolor{bordeaux}{HTML}{5B1A2D}

% --- Page style: plain, centered page numbers (no headers) ------------------
\pagestyle{plain}

% --- First-page footer: license line beneath the centered page number -------
\fancypagestyle{titlepage}{%
    \fancyhf{}%
    \renewcommand{\headrulewidth}{0pt}%
    \renewcommand{\footrulewidth}{0pt}%
    \fancyfoot[C]{%
        \thepage\\[0.4em]
        {\footnotesize\itshape\color{bordeaux} Vinotheca \quad\textbullet\quad published under CC BY-NC 4.0}%
    }%
}

% --- Section formatting (TasteRank style + bordeaux) ------------------------
% Section: small caps with bordeaux rule beneath, lettered prefix manually
\titleformat{\section}[block]
    {\normalsize\scshape\color{bordeaux}}{\thesection}{0.6em}{}%
    [\vspace{2pt}{\color{bordeaux}\hrule}]
\titlespacing*{\section}{0pt}{1.6em}{0.7em}

\titleformat{\subsection}{\normalsize\bfseries\color{bordeaux}}{\thesubsection}{0.6em}{}
\titlespacing*{\subsection}{0pt}{1.1em}{0.4em}

\titleformat{\subsubsection}{\normalsize\bfseries\itshape\color{bordeaux}}{\thesubsubsection}{0.5em}{}
\titlespacing*{\subsubsection}{0pt}{0.9em}{0.3em}

% --- Use letter-prefixed section numbering (A, B, C, ...) -------------------
\renewcommand{\thesection}{\Alph{section}.}
\renewcommand{\thesubsection}{\Alph{section}.\arabic{subsection}}
\renewcommand{\thesubsubsection}{\Alph{section}.\arabic{subsection}.\arabic{subsubsection}}

% --- Bold vectors/matrices shorthand ----------------------------------------
\newcommand{\vc}[1]{\mathbf{#1}}

\begin{document}
\thispagestyle{titlepage}

% ============================================================================
% TITLE BLOCK (TasteRank style)
% ============================================================================
\begin{center}
    {\color{bordeaux}\rule{\textwidth}{0.4pt}}\\[0.6em]
    {\Large\scshape\color{bordeaux} Bayesian Shrinkage}\\[0.4em]
    {\itshape\color{bordeaux} The Collection --- Rating Adjustment}\\[0.4em]
    Hierarchical Model, Derivation, and Application to Codex Data\\[0.6em]
    Jure Skarabot \quad\textbullet\quad MMXXVI\\[0.4em]
    {\color{bordeaux}\rule{\textwidth}{0.4pt}}
\end{center}
\vspace{0.5em}

% ============================================================================
\section{Setting and Notation}

The codex contains $118$ wines partitioned into $K$ groups (countries, or
analogously grape categories). Let group $i$ have $n_i$ wines, each an
independent rating drawn from the group's underlying distribution, with
sample mean
\[
    \bar{x}_i \;=\; \frac{1}{n_i}\sum_{j=1}^{n_i} x_{ij}.
\]

We want an estimator of each group's true mean rating. The naïve choice is
the sample mean itself: the maximum likelihood estimator under independent
sampling, which is unbiased. But when some $n_i$ are small, the
sample-mean estimator has large variance --- a country with $n = 1$ has
variance equal to the full population variance --- and the resulting
rankings are dominated by sampling noise rather than by real differences
between groups.

Bayesian shrinkage resolves this by imposing a hierarchical model: the
group means are themselves drawn from a common distribution, and each
group's sample mean is combined with that population-level information to
produce a posterior estimate. The result is a point estimator that is
biased but has smaller mean squared error than the MLE across every
realistic configuration of the true means.

\subsection{The Normal--Normal Hierarchical Model}

Assume:
\begin{align*}
    x_{ij} \mid \theta_i &\sim \mathcal{N}(\theta_i,\, \sigma^2),
        \quad \text{independent across } i,\,j \\
    \theta_i &\sim \mathcal{N}(\mu,\, \tau^2),
        \quad \text{independent across } i.
\end{align*}

Here the within-group variance (how much ratings vary within a single
country) is denoted by one parameter, the between-group variance (how
much true country means vary) by another, and the grand mean by a third.
This is the simplest non-trivial hierarchical model and suffices for
rating data bounded roughly in $[8.9,\, 10.0]$, where normality is a
reasonable working assumption.

Given the model, the group's sample mean is a sufficient statistic for
the true mean, and its distribution conditional on the true mean is
\[
    \bar{x}_i \mid \theta_i \;\sim\; \mathcal{N}\!\left(\theta_i,\; \sigma^2 / n_i\right).
\]

% ============================================================================
\section{Posterior Derivation}

\subsection{Conjugacy and the Posterior Mean}

Because the Normal family is self-conjugate, the posterior for the true
mean is also Normal. Applying Bayes' rule with a Normal prior and a
Normal likelihood, the posterior precision is the sum of prior and
likelihood precisions:
\[
    \frac{1}{\operatorname{Var}(\theta_i \mid \bar{x}_i)}
        \;=\; \frac{1}{\tau^2} + \frac{n_i}{\sigma^2}.
\]

The posterior mean is the precision-weighted average of the prior mean
and the data mean:
\[
    \mathbb{E}[\theta_i \mid \bar{x}_i]
        \;=\; \frac{\mu / \tau^2 + n_i\,\bar{x}_i / \sigma^2}{1/\tau^2 + n_i / \sigma^2}.
\]

Multiplying numerator and denominator by the within-group variance and
rearranging:
\[
    \hat{\theta}_i \;=\; \frac{n_i\,\bar{x}_i + (\sigma^2 / \tau^2)\,\mu}{n_i + \sigma^2 / \tau^2}.
\]

Define the shrinkage constant as the ratio of within- to between-group
variance:
\[
    m \;=\; \sigma^2 / \tau^2.
\]

The estimator takes the clean closed form
\[
    \boxed{\;\hat{\theta}_i \;=\; \frac{n_i\,\bar{x}_i + m\,\mu}{n_i + m}\;}.
\]

This is the formula implemented on the site. It is the posterior mean
under the Normal--Normal hierarchical model; equivalently, it is the
Bayes estimator under squared-error loss. As the sample size grows, the
estimator converges to the raw mean; as it vanishes, the estimator
converges to the prior mean. The ratio $m$ has an intuitive reading: it
is the number of observations whose evidence is worth, in aggregate, the
same as the prior --- a prior ``sample size.''

\subsection{Equivalent Weighted-Average Form}

Define
\[
    \alpha_i \;=\; \frac{n_i}{n_i + m}.
\]
Then the estimator becomes
\[
    \hat{\theta}_i \;=\; \alpha_i\,\bar{x}_i + (1 - \alpha_i)\,\mu.
\]

This makes the shrinkage explicit: the estimate is a convex combination
of the group mean and the grand mean, with the weight on the group mean
growing with sample size. The factor $1 - \alpha_i$ is the
\emph{shrinkage fraction}. When the sample size equals the shrinkage
constant $m$, the estimate sits exactly halfway between the group mean
and the grand mean.

% ============================================================================
\section{Empirical Bayes Estimation from Codex Data}

\subsection{Variance Components by Method of Moments}

The formula requires both variance components. In a fully Bayesian
treatment these would have their own priors; the \emph{empirical Bayes}
approach estimates them from the data by method of moments.

With the current $118$ wines:
\begin{align*}
    \hat{\mu}      &= 9.516, \\
    \hat{\sigma}^2 &= \mathrm{MSW} = 0.0616, \\
    \hat{\tau}^2   &= \max\!\bigl(0,\; (\mathrm{MSB} - \mathrm{MSW}) / \bar{n}\bigr) = 0.0039,
\end{align*}

where $\mathrm{MSW}$ is the mean squared error within countries,
$\mathrm{MSB}$ is the mean squared error between country means (weighted
by sample size), and the average sample size across countries with at
least two wines is $10.45$. Singletons are excluded from the
between-group variance estimation because their contribution to the
between-group sum of squares is indistinguishable from within-group
noise.

The empirical-Bayes shrinkage constant is
\[
    \hat{m}_{\text{EB}} \;=\; \hat{\sigma}^2 / \hat{\tau}^2 \;\approx\; 15.69.
\]

\subsection{Choice of a Fixed Shrinkage Constant}

The site uses a fixed $m = 5$ rather than the empirical-Bayes estimate.
With the empirical-Bayes estimate substantially larger than $5$, it
would shrink ranks more aggressively than the site currently does ---
reflecting that within-country variance greatly exceeds between-country
variance, so most of the observed spread in raw country means is
sampling noise rather than real differences in true group means. Two
practical considerations argue against simply adopting the
empirical-Bayes value.

\textbf{Stability of the estimate.} The between-group variance is itself
estimated from only $11$ groups with at least two wines. Its sampling
distribution is wide enough that the empirical-Bayes estimate could
easily double or halve as the codex grows by a dozen wines. A fixed $m$
produces a stable ranking that is not sensitive to small changes in the
data.

\textbf{Visual legibility.} At an empirical-Bayes value near $15.7$,
small-$n$ countries are shrunk so close to the grand mean that the
chart becomes visually flat and conveys little. $m = 5$ preserves more
of the raw signal while still correcting for the most obvious
small-sample distortions. This is a deliberate bias-for-legibility
trade-off.

% ============================================================================
\section{Worked Example}

Applying the estimator with $m = 5$ and grand mean $9.516$ to the
current country data:

\begin{center}
\begin{tabular}{l r r r r r}
\toprule
Country ($i$) & $n_i$ & $\bar{x}_i$ & $\alpha_i$ & $\hat{\theta}_i$ & shift \\
\midrule
France       & $32$ & $9.591$ & $0.865$ & $9.581$ & $-0.010$ \\
Italy        & $26$ & $9.485$ & $0.839$ & $9.490$ & $+0.005$ \\
USA          & $25$ & $9.548$ & $0.833$ & $9.543$ & $-0.005$ \\
Germany      &  $8$ & $9.537$ & $0.615$ & $9.529$ & $-0.008$ \\
Austria      &  $5$ & $9.480$ & $0.500$ & $9.498$ & $+0.018$ \\
Spain        &  $5$ & $9.280$ & $0.500$ & $9.398$ & $+0.118$ \\
South Africa &  $3$ & $9.333$ & $0.375$ & $9.448$ & $+0.114$ \\
Argentina    &  $3$ & $9.767$ & $0.375$ & $9.610$ & $-0.157$ \\
New Zealand  &  $3$ & $9.333$ & $0.375$ & $9.448$ & $+0.114$ \\
Australia    &  $3$ & $9.433$ & $0.375$ & $9.485$ & $+0.052$ \\
Chile        &  $2$ & $9.750$ & $0.286$ & $9.583$ & $-0.167$ \\
Slovenia     &  $1$ & $8.900$ & $0.167$ & $9.413$ & $+0.513$ \\
Greece       &  $1$ & $9.200$ & $0.167$ & $9.463$ & $+0.263$ \\
Croatia      &  $1$ & $9.400$ & $0.167$ & $9.497$ & $+0.097$ \\
\bottomrule
\end{tabular}
\end{center}

\vspace{0.5em}
The $\alpha$ column is the effective weight on each country's own data.
France ($n = 32$) has $\alpha$ close to $0.86$, so its estimate barely
moves. Countries with $n = 1$ have $\alpha$ close to $0.17$, meaning
about $83\%$ of their estimate comes from the grand mean --- a direct
reflection of how little we know about them. Notice that the shrinkage
magnitude depends on both the sample size \emph{and} how far the raw
mean sits from the grand mean: Argentina's three wines average far above
the grand mean, so it gets pulled visibly toward it; Austria sits near
the grand mean and barely moves despite a similar sample size.

% ============================================================================
\section{Properties of the Estimator}

Under the assumed model, the estimator has three notable properties.

\subsection{Bias Toward the Grand Mean}

The bias vanishes as the sample size grows. Explicitly, the expected
value of the estimator minus the true mean equals the product of the
shrinkage fraction and the distance between the true mean and the grand
mean:
\[
    \mathbb{E}[\hat{\theta}_i] - \theta_i \;=\; (1 - \alpha_i)(\mu - \theta_i).
\]

\subsection{Lower Variance than the MLE}

The variance of the estimator conditional on the true mean is the MLE's
variance multiplied by the square of $\alpha$:
\[
    \operatorname{Var}(\hat{\theta}_i \mid \theta_i) \;=\; \alpha_i^2 \cdot \frac{\sigma^2}{n_i}.
\]
Since $\alpha_i$ is strictly less than one, this is a meaningful
reduction.

\subsection{Admissibility}

There exists no estimator that dominates the posterior mean in mean
squared error uniformly over all possible configurations of the true
group means. This is a standard decision-theoretic result in the Normal
means problem.

\subsection{Implications for Ranking}

For group rankings --- the use case here --- what matters is that the
posterior mean is the Bayes estimator under squared-error loss, and that
it orders groups in a way that accounts for sampling uncertainty. A
country ranked high by the estimator is high either because its true
mean is high, or because its sample size is large enough that the
sample mean is trustworthy; either is a defensible reason to display it
high.

% ============================================================================
\section{Limitations}

\subsection{Bounded and Censored Ratings}

The normality assumption is a working convenience. Ratings are bounded
and right-censored near $10.0$ --- there is no rating above $10.0$, and
my own $9.9$s are de facto truncated draws from whatever unbounded
latent variable they represent. For bounded ratings with potential
ceiling effects, a Beta-Binomial or ordered-logit model would be more
principled but would not change the qualitative ranking meaningfully at
current sample sizes.

\subsection{Exchangeability Across Groups}

The model assumes the true group means are exchangeable across
countries, which is a convenient fiction. A Burgundy and a Mosel wine
are drawn from visibly different underlying distributions. A more
refined model would introduce further hierarchy --- continent, climate
zone, price tier --- at the cost of requiring substantially more data
to identify the additional variance components.

\subsection{Selection in the Underlying Sample}

Finally, the ratings themselves are not independent samples from a
well-defined population. They reflect my purchase patterns, which are
non-random --- I buy what I expect to enjoy --- and my tasting context,
which varies. No amount of statistical machinery fixes that; shrinkage
only addresses the sample-size problem conditional on whatever
selection process generated the data in the first place.

% ============================================================================
\section{References}

\begin{itemize}
    \item Efron, B. \& Morris, C. (1975). ``Data Analysis Using Stein's
        Estimator and Its Generalizations.'' \textit{Journal of the
        American Statistical Association} 70 (350): 311--319.
    \item Gelman, A., Carlin, J.\,B., Stern, H.\,S., Dunson, D.\,B.,
        Vehtari, A., \& Rubin, D.\,B. (2013). \textit{Bayesian Data
        Analysis}, 3rd ed., chapter 5.
    \item Robinson, D. (2017). \textit{Introduction to Empirical Bayes:
        Examples from Baseball Statistics}.
\end{itemize}

\end{document}
